\section{Set Theory and Logic}

Uppercase letters usually denote sets; lowercase letters usually denote their elements.

\begin{notationtable}
$\N=\{1,2,3,\dots\}$ & Positive integers. \\
$\Nzero=\{0,1,2,3,\dots\}$ & Nonnegative integers. \\
$\Z$ & Integers. \\
$\Q$ & Rational numbers. \\
$\R$ & Real numbers. \\
$\C$ & Complex numbers. \\
$\varnothing,\ \emptyset$ & Empty set. \\
$x\in A$ & $x$ is an element of $A$. \\
$x\notin A$ & $x$ is not an element of $A$. \\
$A\subseteq B$ & $A$ is a subset of $B$; equality is allowed. \\
$A\subsetneq B$ & $A$ is a proper subset of $B$. \\
$A\supseteq B$ & $A$ is a superset of $B$; equality is allowed. \\
$A\supsetneq B$ & $A$ is a proper superset of $B$. \\
$A\cup B$ & Union. \\
$A\cap B$ & Intersection. \\
$(A_i)_{i\in I}$ & Indexed family of sets. \\
$\displaystyle\bigcup_{i\in I} A_i$ & Union of an indexed family of sets. \\
$\displaystyle\bigcap_{i\in I} A_i$ & Intersection of an indexed family of sets. \\
$A\setminus B,\ A-B$ & Set difference. \\
$A\triangle B,\ A\mathbin{\Delta}B$ & Symmetric difference. \\
$A^c$ & Complement of $A$ in the ambient set. \\
$\Pow(A),\ 2^A$ & Power set. \\
$A\times B$ & Cartesian product. \\
$A_1\times\cdots\times A_n,\ A^n$ & Finite Cartesian product; $A^n$ is the $n$-fold Cartesian power of $A$. \\
$\displaystyle\prod_{i\in I} A_i$ & Cartesian product of an indexed family of sets. \\
$A\sqcup B,\ A\mathbin{\dot\cup}B$ & Disjoint union. \\
$\displaystyle\bigsqcup_{i\in I} A_i$ & Disjoint union of an indexed family of sets. \\
$A\cap B=\varnothing$ & $A$ and $B$ are disjoint. \\
$\{x\in A\mid P(x)\}$ & Elements $x\in A$ satisfying $P(x)$. \\
\addlinespace[0.35em]
$R\subseteq X\times Y$ & Binary relation between $X$ and $Y$. \\
$xRy$ & $x$ is related to $y$ by the relation $R$. \\
$R^{-1}$ & Inverse relation of $R$. \\
$[x]_R,\ [x],\ R[x]$ & Equivalence class of $x$ under $R$; abbreviated forms are used when the relation is clear. \\
$x\sim y$ & Equivalent elements under an understood equivalence relation. \\
$A/{\sim}$ & Set of equivalence classes. \\
\addlinespace[0.35em]
$\dom f$ & Domain of $f$. \\
$\codom f$ & Codomain of $f$. \\
$f(A)$ & Image of $A$ under $f$. \\
$f^{-1}(B)$ & Preimage of $B$ under $f$; this notation does not require $f$ to be invertible. \\
$\im f$ & Image of the map $f$. \\
$f^{-1}\colon Y\to X$ & Inverse function of a bijection $f\colon X\to Y$. \\
$f\colon X\hookrightarrow Y$ & Injective map. \\
$f\colon X\twoheadrightarrow Y$ & Surjective map. \\
$f\colon X\bijarrow Y$ & Bijective map. \\
$f\colon X\xrightarrow{\sim}Y$ & Isomorphism; for sets, equivalently a bijection. \\
$\indicator{A}(x)$ & Indicator function of $A$. \\
\addlinespace[0.35em]
$\forall x\in A$ & For every $x\in A$. \\
$\exists x\in A$ & There exists $x\in A$. \\
$\exists!x\in A$ & There exists a unique $x\in A$. \\
$\nexists x\in A$ & There does not exist $x\in A$. \\
$P\Rightarrow Q$ & $P$ implies $Q$ in ordinary mathematical writing. \\
$P\Leftarrow Q$ & $Q$ implies $P$ in ordinary mathematical writing. \\
$P\nRightarrow Q$ & $P$ does not imply $Q$. \\
$P\Leftrightarrow Q$ & $P$ if and only if $Q$ in ordinary mathematical writing. \\
$P\land Q$ & Logical conjunction. \\
$P\lor Q$ & Logical disjunction. \\
$\neg P$ & Logical negation. \\
$\top,\ \bot$ & Truth and falsehood constants, respectively, in formal logic. \\
$\varphi\to\psi,\ \varphi\leftrightarrow\psi$ & Conditional and biconditional connectives in formal logic. \\
$\mathcal M\models\varphi,\ \mathcal M\not\models\varphi$ & $\mathcal M$ satisfies, or does not satisfy, the formula $\varphi$. \\
$\Gamma\models\varphi,\ \Gamma\not\models\varphi$ & $\varphi$ is, or is not, a semantic consequence of $\Gamma$. \\
$\Gamma\vdash\varphi,\ \Gamma\nvdash\varphi$ & $\varphi$ is, or is not, syntactically derivable from $\Gamma$. \\
$\lightning$ & Contradiction reached in a proof or argument; distinct from the formal falsehood constant $\bot$. \\
$\therefore$ & Therefore; indicates a conclusion from preceding statements. \\
$\because$ & Because; indicates a reason or justification. \\
\addlinespace[0.35em]
$|A|,\ \#A$ & Cardinality. \\
$\aleph_0$ & Cardinality of a countably infinite set. \\
$\mathfrak{c}$ & Cardinality of the continuum. \\
\end{notationtable}
